\scanpage{129}
sous-module \inserted{\unsure{ess. const.}} d'un \(R^mg_*(\unsure{\mathbb Q_l|Q_l})\),
\(g:V\mathbin{\text{---}}U\) un morphisme \uline{projectif} de dim
\inserted{relative} \(\le d\).
On a donc des suites exactes
\[
0\longrightarrow \underset{\unsure{/}}{G}\mathbin{\text{---}}
R^mg_*(\unsure{\mathbb Q_l|Q_l})\mathbin{\text{---}}H\mathbin{\text{---}}0
\]
\[
0\mathbin{\text{---}}\underset{*}{I}\mathbin{\text{---}}
\underset{*}{G}\mathbin{\text{---}}F\longrightarrow0
\]
et \uline{de} nouvelles applications du lemme 2 nous ramènent au cas des \(I\) et \(G\),
\(R^nf_*(F)\) lorsque \deleted{\(I\subset G\)}
\(F\) \unsure{ess. const.} \(\subset R^mg_*(\unsure{\mathbb Q_l|Q_l})\).
Pour continuer, il faut utiliser le
\uline{``th.'' de complète réductibilité}
pour nous ramener au cas où \(F\) est de la forme \(R^mg_*(\unsure{\mathbb Q_l|Q_l})\)
on doit donc examiner le cas
\[
R^p_!f_*(R^qg_*(Q_l))\qquad p,q\ge0
\]

\deleted{(\unsure{N°|1°}) Cas \(m=0\). Nous procédons par réc. \uline{descendante}
sur \(p+q\) \unsure{supprimé}, en supposant que c'est vrai pour \((p',q')\)
si \(p'>p\) \uline{ou} \(q\)}
% archive foliation: 128
